% Provenance (never printed): see sources/outline_sources.md (Year 2: radiative forcing
% quantified; feedbacks and climate sensitivity; detection and attribution; scenarios; the
% remaining carbon budget; components of sea-level rise; physical science only). Every statement
% from IPCC AR6 WG1 (SPM; chapters 3, 5, 7 and 9; Annex III data), ledger rows in
% sources/ledger/book2.md (ipcc:*) and sources/ledger/book4.md (ar6*:*). Curves and numbers from
% figdata/bachelor-2/anthropogenic.py and the Book 4 second-writer scratch calculator calc/ch27.py.
\chapter{人为气候变化}\label{ch:b2:anthropogenic-climate-change}

自 1970 年以来，全球地表的升温速度快于至少过去两千年中的任何一个五十年期
（\emph{高度可信}）。中学各册已经讲清了发生了什么变化，以及原因。
本章把这一叙述转化为一笔收支账：人为驱动因子每平方米增加多少瓦特，
变暖的星球又向太空送回多少，每吨二氧化碳会带来多少升温，
以及上升的海平面如何在变暖的海水与熔化的冰之间分配。每个数值都附有不确定度，
并采用评估报告中经过校准的措辞表述。
% ledger: ipcc:fastest-2000yr, ipcc:likelihood

\begin{recall}
\textbf{来自中学各册}：辐射强迫（二氧化碳加倍时为 $\qty{3.93}{W.m^{-2}}$）及
其对数定律；气候反馈；\omterm{def:b2:anthropogenic-climate-change:ecs}{平衡气候敏感度}，最佳估计为 \qty{3}{\celsius}，
\emph{可能}在 2.5 至 \qty{4}{\celsius} 之间；\omterm{def:b2:anthropogenic-climate-change:da}{归因}、排放情景与预估；
校准的措辞为：\emph{可能}对应 66 至 \qty{100}{\percent}，\emph{很
可能}对应 90 至 \qty{100}{\percent}；区间 $[x \text{ 至 } y]$ 即为很可能区间。人类
活动使 2010--2019 年的气候相对 1850--1900 年升温了 \qty{1.07}{\celsius}。
\textbf{来自物理学各册}：黑体与斯特藩–玻尔兹曼定律；热膨胀。
\textbf{来自生物学各册}：碳循环、大气分数、碳汇，以及生物圈的
碳收支。\textbf{来自\cref{ch:b2:deep-time-climate}}：慢碳循环，
它只能在数十万年的时间尺度上吸回过量的二氧化碳。
\end{recall}
% ledger: ipcc:erf-2xco2, ipcc:log-co2, ipcc:ecs, ipcc:likelihood, ipcc:attribution

\section{辐射强迫的定量化}\label{sec:b2:anthropogenic-climate-change:forcing}

\begin{definition}[有效辐射强迫]\label{def:b2:anthropogenic-climate-change:erf}
扰动（温室气体、气溶胶或太阳的变化）的\emph{有效辐射强迫}\index{有效辐射强迫}（ERF）
是指在大气、云与陆面完成调整、而全球地表温度尚未响应之前，
大气层顶净向下辐射通量的变化量，单位为 \unit{W.m^{-2}}。
计入的调整是由该因子本身引起的，
与地表温度
无关。
\end{definition}
% ledger: ar6ch7:erf-def

该定义以是否依赖地表温度来区分强迫与响应，而不以变化快慢为准。
二氧化碳的瞬时强迫是气体加入后通量刚刚发生的变化；
随后在数月内降温的平流层使其第一次改变，
对流层、其中的水汽与云则再次改变它。对于二氧化碳加倍的情形，
评估报告给出平流层调整后的强迫为 \qty{3.75}{W.m^{-2}}，
对流层中部分相互抵消的调整：来自其升温的 \qty{-0.60}{W.m^{-2}}，来自水汽的 +0.22，
来自云的 +0.45，来自陆面的 +0.11，合计 \qty{+0.18}{W.m^{-2}}。
ERF 为 $3.75 + 0.18 = \qty{3.93}{W.m^{-2}}$，比上一次评估的
\qty{3.71}{W.m^{-2}} 高出 \qty{6}{\percent}（见\cref{exo:b2:anthropogenic-climate-change:1}）。
% ledger: ar6ch7:2xco2, ar6ch7:erf-def

自 1750 年以来，各因子的强迫均已逐年重建。下图
给出 2019 年的情况：二氧化碳为 \qty{2.16}{W.m^{-2}}，其他温室气体——甲烷、
一氧化二氮、卤代气体与臭氧——合计 \qty{1.69}{W.m^{-2}}，而气溶胶为
\qty{-1.06}{W.m^{-2}}。人为强迫的总和为 \qty{2.72}{W.m^{-2}}，其
很可能区间为 1.96 至 \qty{3.48}{W.m^{-2}}，区间较宽主要源于气溶胶。气溶胶的
ERF 在 1750--2014 年间评估为 $-1.3$ [$-2.0$ 至 $-0.6$] \unit{W.m^{-2}}（\emph{中等可信}），
其中绝大部分通过对云的作用实现，云在气溶胶提供更多云滴时会变亮。
火山喷发会带来短暂的负向尖峰（见\cref{ch:b2:eruptive-styles}）：记录中最深的一次
为 \qty{-4.8}{W.m^{-2}}（1816 年），紧随 1815 年坦博拉火山的喷发。
% ledger: data:b4-ar6-erf, ipcc:erf-total, ipcc:erf-agents, ar6ch7:aerosol, usgs:msh-comparisons

\begin{omfigure}
\centering
\begin{tikzpicture}
\begin{axis}[omrecord, width=0.92\linewidth, height=6.2cm, xmin=1750, xmax=2020, ymin=-5.2,
  ymax=3.8, xlabel={年}, ylabel={自 1750 年以来的 ERF（\unit{W.m^{-2}}）},
  xtick={1750,1800,...,2000}, /pgf/number format/1000 sep={}, xtick pos=bottom, ytick pos=left,
  tick label style={font=\scriptsize}, label style={font=\small},
  legend style={at={(0.5,-0.24)}, anchor=north, font=\scriptsize, draw=black!30,
  legend columns=3, /tikz/every even column/.append style={column sep=6pt}},
  legend cell align=left]
% computed in figdata/bachelor-2/anthropogenic.py part erf; ledger: data:b4-ar6-erf
\addplot[name path=hi, draw=none, forget plot] table[x=year, y=hi] {figdata/out/bachelor-2-anthropogenic-erf.dat};
\addplot[name path=lo, draw=none, forget plot] table[x=year, y=lo] {figdata/out/bachelor-2-anthropogenic-erf.dat};
\addplot[black!15, forget plot] fill between[of=hi and lo];
\addplot[black!45, line width=0.4pt] table[x=year, y=natural] {figdata/out/bachelor-2-anthropogenic-erf.dat};
\addlegendentry{自然因子（火山、太阳）}
\addplot[omThm, line width=1pt] table[x=year, y=co2] {figdata/out/bachelor-2-anthropogenic-erf.dat};
\addlegendentry{二氧化碳}
\addplot[omThm!50!orange, line width=1pt, densely dashed] table[x=year, y=other_ghg] {figdata/out/bachelor-2-anthropogenic-erf.dat};
\addlegendentry{其他温室气体与臭氧}
\addplot[omDef, line width=1pt] table[x=year, y=aerosol] {figdata/out/bachelor-2-anthropogenic-erf.dat};
\addlegendentry{气溶胶}
\addplot[black, line width=1.4pt] table[x=year, y=anthro] {figdata/out/bachelor-2-anthropogenic-erf.dat};
\addlegendentry{人为总强迫（灰色：5--95\,\%）}
\draw[black!40] (axis cs:1750,0) -- (axis cs:2020,0);
\node[anchor=west, font=\scriptsize, text=black!60, inner sep=1pt] at (axis cs:1817.5,-4.75) {$\leftarrow$ 坦博拉};
\end{axis}
\end{tikzpicture}
\par\smallskip{\small 按因子给出的 1750 至 2019 年\omterm{def:b2:anthropogenic-climate-change:erf}{有效辐射强迫}，相对 1750 年。
人为总强迫在 2019 年达到 \qty{2.72}{W.m^{-2}}，阴影表示很可能区间。
数据：IPCC Sixth Assessment Report, Working Group~I, Annex~III extended data (CC BY 4.0).}
\end{omfigure}
% values: anthro 2.7204 (1.9569-3.4819), co2 2.1563, other_ghg 1.6850, aerosol -1.0583, volcanic -4.8162 in 1816 (anthropogenic.py numbers)

\section{能量收支}\label{sec:b2:anthropogenic-climate-change:energy}

正强迫使行星获得能量。随着升温，行星向太空辐射的能量增加，
获得的能量随之减少；只有两者达到平衡时，气候才会稳定下来。

\begin{definition}[地球能量失衡；气候反馈参数]\label{def:b2:anthropogenic-climate-change:eei}
\emph{地球能量失衡}\index{地球能量失衡} $N$ 是指在大气层顶流入气候系统的净能量通量，
经全球平均并取一年或更长时间的平均值，
单位为 \unit{W.m^{-2}}：吸收的太阳辐射减去向外射出的红外线。它是海洋、
陆地、冰和空气获得热量的速率，以每平方米地球表面计。
\end{definition}

\begin{definition}[气候反馈参数]\label{def:b2:anthropogenic-climate-change:lambda}
\emph{气候反馈参数}\index{气候反馈参数} $\lambda$ 是指
全球地表每升温一度时，大气层顶净向下通量的变化量，
由气候系统对该升温的响应引起，单位为 \unit{W.m^{-2}.\celsius^{-1}}。
稳定的气候满足 $\lambda < 0$：升温会增加能量向太空的损失。
\end{definition}

高中分册以 $\lambda$ 的大小（一个正数）写平衡关系；从此处起，
保留符号，与评估报告的写法一致，后者用参数 $\alpha$ 表示。
% ledger: ar6ch7:energy-budget

\begin{proposition}[全球能量平衡]\label{prop:b2:anthropogenic-climate-change:balance}
设 $\Delta F$ 为\omterm{def:b2:anthropogenic-climate-change:erf}{有效辐射强迫}，$\Delta T$ 为
相对于平衡态的全球地表升温。若射出通量的响应关于 $\Delta T$ 为线性，
\[
  \Delta N = \Delta F + \lambda\,\Delta T .
\]
在新的平衡态下 $\Delta N = 0$，因此 $\Delta T_{\text{eq}} = -\Delta F/\lambda$。
\end{proposition}
\begin{proof}
由定义，在地表温度固定时，强迫使净向下通量改变 $\Delta F$，
而升温引起的响应到一阶近似下使其改变 $\lambda\,\Delta T$；
失衡量即二者之和。令其为零，即得平衡态。
\end{proof}

测量 $N$ 即测量系统所获得的热量，而海洋几乎储存了全部热量。
1971 至 2018 年间，气候系统共获得 \qty{435}{ZJ}（泽焦耳，\qty{e21}{J}）：海洋
占 \qty{91}{\percent}，其中上层 \qty{700}{m} 内为 \qty{242}{ZJ}，以下为 \qty{154}{ZJ}；陆地
占 \qty{5}{\percent}，熔化的冰占 \qty{3}{\percent}，大气占 \qty{1}{\percent}。平均
失衡在 1971--2018 年间为 \qty{0.57}{W.m^{-2}}，且有所增大：2006--2018 年为 \qty{0.79}{W.m^{-2}}
[0.52 至 1.06]（见\cref{exo:b2:anthropogenic-climate-change:2}）。
% ledger: ar6ch7:inventory, ipcc:heating

\begin{omfigure}
\centering
\begin{tikzpicture}[font=\small, >=Stealth]
\draw[fill=omDef!8, draw=omDef!60] (-3.2,0) rectangle (3.2,1.6);
\node[align=center] at (0,1.05) {气候系统};
\node[align=center, font=\scriptsize] at (0,0.5) {海洋 \qty{91}{\percent}，陆地 \qty{5}{\percent}，冰 \qty{3}{\percent}，大气 \qty{1}{\percent}};
\draw[->, line width=2.2pt, omThm] (-2.0,3.3) -- node[pos=0.12, left=3pt, align=right, font=\scriptsize] {强迫 $\Delta F$} (-2.0,1.7);
\draw[->, line width=2.2pt, omDef] (2.0,1.7) -- node[pos=0.88, right=3pt, align=left, font=\scriptsize] {响应 $\lambda\,\Delta T$} (2.0,3.3);
\draw[->, line width=1.4pt, black!70] (0,-0.1) -- node[right=3pt, align=left, font=\scriptsize] {储存：$\Delta N = \Delta F + \lambda\,\Delta T$} (0,-1.0);
\draw[black!40, dashed] (-3.4,2.4) -- (3.4,2.4) node[right, font=\scriptsize, text=black!60] {大气层顶};
\end{tikzpicture}
\par\smallskip{\small 受强迫气候的能量收支（示意图）。强迫使
大气层顶增加能量，变暖的行星的响应移除一部分（$\lambda\,\Delta T$ 为
负值），二者之差被储存，主要在海洋中。图中为 1971
至 2018 年间所获热量的占比。}
\end{omfigure}
% ledger: ar6ch7:inventory

\begin{method}[由历史能量收支求气候敏感度]\label{met:b2:anthropogenic-climate-change:energy-budget}
\begin{enumerate}
  \item 选取两个时段：一个接近平衡态的参考时段，一个近期时段；取
        强迫变化 $\Delta F$、失衡变化 $\Delta N$ 及全球温度变化 $\Delta T$。
  \item 由能量平衡式解出反馈参数：$\lambda = (\Delta N - \Delta F)/\Delta T$。
  \item 于是倍增敏感度为 $-F_{2\times}/\lambda$；
        强迫发生时的升温按 $F_{2\times}$ 折算后，为 $F_{2\times}\Delta T/\Delta F$。
\end{enumerate}
\end{method}

自 1850--1900 年至 2006--2019 年，评估报告取 $\Delta N = \qty{0.59}{W.m^{-2}}$，$\Delta T =
\qty{1.03}{\celsius}$，$\Delta F = \qty{2.20}{W.m^{-2}}$。于是 $\lambda = (0.59 - 2.20)/1.03 =
\qty{-1.56}{W.m^{-2}.\celsius^{-1}}$，敏感度为 $3.93/1.56 = \qty{2.5}{\celsius}$，
与评估报告得出的数值一致，其很可能区间为 1.6 至 \qty{4.8}{\celsius}。但这是一个
\emph{有效}敏感度，针对迄今为止升温所对应的反馈。升温
最强之处恰为反馈最负的区域，即热带西太平洋，升温较弱的则在
东太平洋和南大洋上；随着这一格局演变，反馈参数预计会变得不那么负（\emph{高信度}），
因此长期敏感度高于
由历史估计值（见\cref{exo:b2:anthropogenic-climate-change:7}）。
% ledger: ar6ch7:energy-budget, ar6ch7:pattern, ipcc:erf-2xco2
% values: lambda_energy -1.563, ecs_energy 2.514 (anthropogenic.py numbers)

\section{反馈与敏感度}\label{sec:b2:anthropogenic-climate-change:feedbacks}

\begin{definition}[普朗克响应；水汽反馈、递减率反馈与云反馈]\label{def:b2:anthropogenic-climate-change:feedbacks}
反馈参数被分解为随地表温度变化的各量的贡献。
\emph{普朗克响应}\index{普朗克响应}是指地表与整个对流层以相同幅度升温时，
向太空额外发射的红外线。
\emph{水汽反馈}\index{水汽反馈}是指由大气水汽随温度增加而引起的
通量变化。\emph{递减率反馈}\index{递减率反馈}是指升温在高度上不均匀时
所引起的通量变化。\emph{云反馈}\index{云反馈}是指
由云量、云高及
云的亮度变化所引起的通量变化。
\end{definition}

\begin{proposition}[普朗克响应]\label{prop:b2:anthropogenic-climate-change:planck}
若地球以固定比例 $1 - g$ 向太空发射其地表的红外线 $\sigma T^4$，地表温度为
$T$ 时，其\omterm{def:b2:anthropogenic-climate-change:feedbacks}{普朗克响应}为
\[
  \lambda_{\text P} = -(1 - g)\,4\sigma T^3 .
\]
取 $g = 0.4$，$T = \qty{288}{K}$，则 $\lambda_{\text P} = \qty{-3.25}{W.m^{-2}.K^{-1}}$。
\end{proposition}
\begin{proof}
射出通量为 $(1 - g)\sigma T^4$；其关于 $T$ 的导数为 $(1 - g)\,4\sigma
T^3$，额外的损失对净向下通量的贡献为负。数值上 $0.6
\times 4 \times \num{5.670e-8} \times 288^3 = 3.25$。
\end{proof}
% ledger: ar6ch7:planck, const:sigma, ipcc:budget-surface
% values: planck -3.2509 (anthropogenic.py numbers)

分数 $g$ 即归一化的温室效应：大气保留了 \qty{159}{W.m^{-2}}，
占地表所发出的 \qty{398}{W.m^{-2}}，即 $g \approx 0.4$。评估报告采用详细的
辐射模式，得到 $-3.22$ [$-3.4$ 至 $-3.0$]（\emph{高信度}）。水汽遵循
在全球尺度上服从克劳修斯–克拉珀龙定律，水汽增多则吸收更多红外线：
\omterm{def:b2:anthropogenic-climate-change:feedbacks}{水汽反馈}为强正值，由卫星测得 \qty{1.85}{W.m^{-2}.\celsius^{-1}}。
热带地区，对流层上层的增温大于地表，地表随之辐射得更多：
\omterm{def:b2:anthropogenic-climate-change:feedbacks}{递减率反馈}在热带为负，在极地附近为正，全球平均约为 $-0.50$。两者
两者合并评估，为 \num{1.30} [1.1 至 1.5]，因为一个模式在其中一项中所归入的，
往往会从另一项中扣除。熔化的雪和冰使地表变暗，贡献 +0.35；云的贡献为 \num{0.42}
[$-0.10$ 至 $0.94$]，为最大的不确定性。\omterm{def:b2:anthropogenic-climate-change:feedbacks}{云反馈}净值为负的情形\emph{极不可能}出现。
% ledger: ar6ch7:planck, ar6ch7:wv-lr, ar6ch7:table710, ipcc:cloud-negative, ipcc:budget-surface

\begin{omfigure}
\centering
\begin{tikzpicture}
\begin{axis}[xbar, width=0.78\linewidth, height=5.4cm, xmin=-3.7, xmax=1.8, ymin=0.4, ymax=6.6,
  ytick={1,...,6}, yticklabels={净值,生物地球物理与化学,云,地表反照率,
  水汽 + 递减率,普朗克}, y tick label style={font=\scriptsize},
  tick label style={font=\scriptsize}, xlabel={反馈参数（\unit{W.m^{-2}.\celsius^{-1}}）},
  label style={font=\small}, bar width=7pt, axis x line=bottom, axis y line=left,
  axis line style={-}, xmajorgrids, grid style={gray!22, line width=0.3pt}, extra x ticks={0},
  extra x tick style={grid style={black!60, line width=0.5pt}}]
% ledger: ar6ch7:table710 (central values and very likely ranges)
\addplot[fill=omDef!55, draw=omDef!80!black, error bars/.cd, x dir=both, x explicit,
  error bar style={black!70}] coordinates {
  (-3.22,6) += (0.22,0) -= (0.18,0)
  (-0.01,2) += (0.26,0) -= (0.26,0)
  (-1.16,1) += (0.65,0) -= (0.65,0)};
\addplot[fill=omThm!55, draw=omThm!80!black, error bars/.cd, x dir=both, x explicit,
  error bar style={black!70}] coordinates {
  (1.30,5) += (0.20,0) -= (0.20,0)
  (0.35,4) += (0.25,0) -= (0.25,0)
  (0.42,3) += (0.52,0) -= (0.52,0)};
\end{axis}
\end{tikzpicture}
\par\smallskip{\small 由 IPCC 评估的气候反馈（第六次评估报告，第一工作组，
表 7.10），中心值与很可能区间。正反馈放大变暖
（红色），负反馈则抑制变暖（蓝色）；净值为各项之和。}
\end{omfigure}

\begin{definition}[平衡气候敏感度]\label{def:b2:anthropogenic-climate-change:ecs}
称\emph{平衡气候敏感度}\index{平衡气候敏感度}（ECS）为
在新的平衡状态下，大气中二氧化碳含量相对工业化前水平加倍
所引起的全球地表增温，且假定冰盖保持不变。
\end{definition}

\begin{proposition}[敏感度与反馈]\label{prop:b2:anthropogenic-climate-change:ecs}
若反馈参数为常数，则 $\text{ECS} = -F_{2\times}/\lambda$。取 $F_{2\times} =
\qty{3.93}{W.m^{-2}}$，上述各项反馈的净值为 $\lambda = \qty{-1.16}{W.m^{-2}.\celsius^{-1}}$，
则 $\text{ECS} = \qty{3.4}{\celsius}$；仅\omterm{def:b2:anthropogenic-climate-change:feedbacks}{普朗克响应}即给出 $3.93/3.22 =
\qty{1.2}{\celsius}$。
\end{proposition}
\begin{proof}
在平衡态下应用\cref{prop:b2:anthropogenic-climate-change:balance}，取 $\Delta F =
F_{2\times}$。
\end{proof}
% ledger: ar6ch7:table710, ipcc:erf-2xco2
% values: ecs_process 3.388 (anthropogenic.py numbers)

反馈只是证据之一；1850 年以来的增温、能量收支与
过去的气候是其他证据。综合起来，评估的 ECS 为 \qty{3}{\celsius}，
\emph{可能}为 2.5 至 \qty{4}{\celsius}（\emph{高信度}），\emph{很可能}为 2 至 \qty{5}{\celsius}。
由于 ECS 与 $1/\lambda$ 成正比，对称的 $\lambda$ 不确定性
会使 ECS 的范围向高值一侧拉伸（见\cref{exo:b2:anthropogenic-climate-change:12}）。
% ledger: ipcc:ecs

\begin{definition}[瞬态气候响应]\label{def:b2:anthropogenic-climate-change:tcr}
称\emph{瞬态气候响应}\index{瞬态气候响应}（TCR）为
大气中二氧化碳含量加倍时刻的全球地表增温，此时二氧化碳
自工业化前水平起每年增加 \qty{1}{\percent}，约经过七十年。
\end{definition}

TCR 小于 ECS，因为在加倍达到时海洋仍在吸收
热量。其最佳估计为 \qty{1.8}{\celsius}，\emph{可能}为 1.4 至 \qty{2.2}{\celsius}
（\emph{高信度}）；仅历史能量收支即给出 $3.93 \times 1.03/2.20 =
\qty{1.8}{\celsius}$，结果亦同。
% ledger: ar6ch7:tcr, ar6ch7:energy-budget
% values: tcr_energy 1.840 (anthropogenic.py numbers)

\section{检测与归因}\label{sec:b2:anthropogenic-climate-change:attribution}

\begin{definition}[检测；归因；气候指纹]\label{def:b2:anthropogenic-climate-change:da}
\emph{检测}\index{检测（气候变化）}是指表明
气候在某种确定的统计意义上已经发生变化，超出其内部变率所能产生的
范围，但不给出原因。\emph{归因}\index{归因}是指评估
对已检测变化中若干原因的贡献，并赋予相应的可能性。\emph{气候
指纹}\index{气候指纹}是指某一驱动因子响应的时空模式，
由气候模式模拟得到，据此在观测中识别出该驱动因子。
\end{definition}

不同驱动因子留下不同的指纹。温室气体增多使对流层增温，
平流层降温，因为更厚的红外“毯子”使高层向太空辐射得更多；
太阳变亮则两者都会增温。观测结果显示第一种模式：对流层已经增温，
至少自 1950 年代以来，平流层降温是\emph{几乎确定}的。
人为温室气体\emph{很可能}是自 1979 年以来对流层增温的主要驱动因子，
而臭氧层耗竭\emph{极可能}是
1979 年至 1990 年代中期平流层下部降温的主要驱动因子。
% ledger: ar6ch3:upper-air

\begin{method}[最优指纹法]\label{met:b2:anthropogenic-climate-change:fingerprinting}
\begin{enumerate}
  \item 分别模拟各驱动因子单独作用下的响应：指纹 $x_1, \dots, x_k$，
        即空间与时间上取值的矢量。
  \item 将观测写为 $y = \sum_i \beta_i x_i + \varepsilon$，其中 $\varepsilon$ 为
        内部变率，其统计量取自长时间的无外强迫模拟。
  \item 用最小二乘法估计标度因子 $\beta_i$，权重取变率的倒数，
        使波动最大的分量权重最小，并给出各自的置信区间。
  \item 若 $\beta_i$ 显著大于零，则表示对驱动因子 $i$ 的响应
        已被检测到；若区间包含 1 而不含 0，则表示模拟响应
        与观测相符，该变化可部分\omterm{def:b2:anthropogenic-climate-change:da}{归因}于该驱动因子。
\end{enumerate}
\end{method}
% ledger: ar6ch3:fingerprint

将该方法用于地表温度，得到开篇所述的\omterm{def:b2:anthropogenic-climate-change:da}{归因}结果：在
1850--1900 年与 2010--2019 年之间观测到的 \qty{1.06}{\celsius} 增温中，人类影响\emph{可能}
导致 0.8 至 \qty{1.3}{\celsius}，最佳估计为 \qty{1.07}{\celsius}；自然驱动因子和
内部变率的贡献，无论正负，均不超过零点几度
（见\cref{exo:b2:anthropogenic-climate-change:9,exo:b2:anthropogenic-climate-change:10}）。
% ledger: ipcc:attribution, ipcc:observed-2010-2019

\section{排放情景与累积排放}\label{sec:b2:anthropogenic-climate-change:scenarios}

\begin{definition}[共享社会经济路径]\label{def:b2:anthropogenic-climate-change:ssp}
称\emph{共享社会经济路径}\index{共享社会经济路径}（SSP）为
人口、经济和技术在本世纪可能发展的五种叙事之一。评估所用的排放情景记为
SSP$x$-$y$：$x$ 为路径编号，$y$ 为该情景在 2100 年达到的近似辐射强迫，
单位为 \unit{W.m^{-2}}。
\end{definition}
% ledger: ar6spm:ssp

五种主要情景从 SSP1-1.9 延伸到 SSP5-8.5。自然科学评估将这些情景视为给定条件，
既不权衡其可能性，也不评价其叙事。对每种情景，气候模式给出
2081--2100 年相对 1850--1900 年的预估增温：SSP1-1.9 下为 \qty{1.4}{\celsius}（很可能在 1.0 至
1.8 之间），SSP5-8.5 下为 \qty{4.4}{\celsius}（3.3 至 5.7）。这些预估
遵循一个简单规律。
% ledger: ar6spm:ssp, ipcc:proj-temp

\begin{definition}[累积 CO$_2$ 排放的瞬态气候响应]\label{def:b2:anthropogenic-climate-change:tcre}
称\emph{累积 \ce{CO2} 排放的瞬态气候响应}\index{累积 CO$_2$ 排放的瞬态气候响应}
（TCRE）是指自工业化前以来累计排放的每单位二氧化碳所对应的全球地表增温，
单位为 \unit{\celsius}，即每
\qty{1000}{Gt} \ce{CO2} 的排放量。
\end{definition}

\begin{proposition}[增温遵循累积排放]\label{prop:b2:anthropogenic-climate-change:tcre}
在良好近似下，二氧化碳引起的增温与排放总量成正比，
与排放路径无关：$\Delta T \approx \text{TCRE} \times E$，其中 $E$ 为
累积排放量。
\end{proposition}
\admitted
% ledger: ar6spm:tcre, ar6ch5:tcre-linear

这一规律并不显然。辐射强迫仅随浓度的对数增长，因此
每增加一吨排放的增温逐渐减小；但随着排放不断累积，海洋吸收热量与碳的效率逐渐降低，
更多排放留在大气中，其增温的延迟效应也随之减弱。两种效应几乎相互抵消，
所有研究都认为这正是该规律的成因。两者的关系近似线性（\emph{高信度}），
每 \qty{1000}{Gt} 二氧化碳\emph{很可能}引起 0.27 至 \qty{0.63}{\celsius} 增温，最佳估计为 \qty{0.45}{\celsius}，或
每 \qty{1000}{PgC} 碳对应 \qty{1.65}{\celsius} 增温
（见\cref{exo:b2:anthropogenic-climate-change:4}）。1850 至 2019 年间
共排放 $\num{2390} \pm \qty{240}{Gt}$ 二氧化碳。
% ledger: ar6ch5:tcre-linear, ar6spm:tcre, ar6spm:budget

\begin{definition}[剩余碳预算]\label{def:b2:anthropogenic-climate-change:budget}
对某一温度限值而言，\emph{剩余碳预算}\index{剩余碳预算}是指
从某一给定日期起仍可排放的累积二氧化碳量，且在给定概率下
使全球增温保持在该限值以下，并已计入
其他人为驱动因素所造成的增温。
\end{definition}

\begin{method}[由 TCRE 求剩余碳预算]\label{met:b2:anthropogenic-climate-change:budget}
\begin{enumerate}
  \item 将已经达到的人为增温（2010--2019 年为 \qty{1.07}{\celsius}）从
        限值中减去。
  \item 减去其他驱动因素在二氧化碳排放
        达到净零之前还将增加的增温。
  \item 除以 TCRE：所得即为自 2010--2019 年中期起的预算。
  \item 减去此后的排放量，以及为模式未包含的
        缓慢地球系统反馈预留的余量（其中永久冻土融化最大）。
\end{enumerate}
\end{method}

评估报告以这种方式综合五个分量：TCRE；历史增温；
其他驱动因素带来的附加增温约为 0.1 至 \qty{0.2}{\celsius}；排放停止后可能继续出现的增温，
平均评估为零（\emph{很可能}在 $\pm
\qty{0.19}{\celsius}$ 之内）；以及地球系统反馈，约为每
摄氏度 \qty{26}{Gt} 二氧化碳。对于 \qty{1.5}{\celsius} 的限值（相对 2010--2019 年高出 \qty{0.43}{\celsius}），
自 2020 年初起的剩余预算在 \qty{50}{\percent} 概率下为 \qty{500}{Gt} 二氧化碳，
对 \qty{2}{\celsius} 则为 \qty{1350}{Gt}。其他驱动因素的取值选择会使这些数值变动
\qty{220}{Gt} 或更多；周末问题将重新推导第一个数值。
% ledger: ar6ch5:budget-parts, ar6spm:budget

\begin{dated}{2025}
% ledger: data:co2-mlo, data:gistemp (values read from the datasets as fetched 2026-10-22; offset computed)
2025 年莫纳罗亚观测站的年平均二氧化碳浓度达到 \qty{427.35}{ppm}，
按对数规律计算，其强迫为 $3.93\log_2(427.35/278) = \qty{2.44}{W.m^{-2}}$。在 NASA 温度记录中，
2024 年为最暖的一年，比 1880--1900 年的平均值高 \qty{1.50}{\celsius}，
这一基准期为该记录中最早的几年；2025 年则为 \qty{1.40}{\celsius}。
\end{dated}

\begin{omfigure}
\centering
\begin{tikzpicture}
\begin{axis}[omrecord, width=0.86\linewidth, height=5.8cm, xmin=0, xmax=4000, ymin=0, ymax=2.6,
  xlabel={自 1850 年以来的累积 \ce{CO2} 排放量（\unit{Gt}）},
  ylabel={相对 1850--1900 年的增温（\unit{\celsius}）}, xtick={0,500,...,4000},
  /pgf/number format/1000 sep={\,}, xtick pos=bottom, ytick pos=left, clip=false,
  tick label style={font=\scriptsize}, label style={font=\small}]
% computed in figdata/bachelor-2/anthropogenic.py part tcre; ledger: ar6spm:tcre, ar6spm:budget
\addplot[name path=hi, draw=none, forget plot] table[x=E_GtCO2, y=dT_hi] {figdata/out/bachelor-2-anthropogenic-tcre.dat};
\addplot[name path=lo, draw=none, forget plot] table[x=E_GtCO2, y=dT_lo] {figdata/out/bachelor-2-anthropogenic-tcre.dat};
\addplot[omThm!18, forget plot] fill between[of=hi and lo];
\addplot[omThm, line width=1.3pt] table[x=E_GtCO2, y=dT_mid] {figdata/out/bachelor-2-anthropogenic-tcre.dat};
\draw[densely dashed, black!60] (axis cs:0,1.5) -- (axis cs:4000,1.5);
\draw[densely dashed, black!60] (axis cs:0,2.0) -- (axis cs:4000,2.0);
\node[anchor=north west, font=\scriptsize, text=black!70] at (axis cs:0,1.5) {\qty{1.5}{\celsius}};
\node[anchor=north west, font=\scriptsize, text=black!70] at (axis cs:0,2.0) {\qty{2}{\celsius}};
\fill[black] (axis cs:2390,1.07) circle (2.2pt);
\draw[black, line width=0.8pt] (axis cs:2150,1.07) -- (axis cs:2630,1.07);
\draw[black, line width=0.8pt] (axis cs:2390,0.8) -- (axis cs:2390,1.3);
\node[anchor=north west, font=\scriptsize, align=left] at (axis cs:2440,1.02) {1850--2019:\\\qty{2390}{Gt}, \qty{1.07}{\celsius}};
\draw[<->, omDef, thick] (axis cs:2390,1.62) -- (axis cs:2890,1.62);
\node[anchor=south, font=\scriptsize, text=omDef, align=center] at (axis cs:2640,1.64) {AR6: \qty{500}{Gt}};
\node[anchor=north west, font=\scriptsize, text=omThm, align=left] at (axis cs:150,2.5) {可能范围（TCRE）：\\每 \qty{1000}{Gt} 对应 0.27 至 \qty{0.63}{\celsius}};
\end{axis}
\end{tikzpicture}
\par\smallskip{\small 增温随累积二氧化碳排放的变化，取评估的 TCRE，
\qty{0.45}{\celsius} 每 \qty{1000}{Gt}（实线）及其可能范围（阴影），并标出人为增温
与迄今排放（圆点，可能范围以条形表示）。箭头标出
自 2020 年起、在 \qty{50}{\percent} 概率下达到 \qty{1.5}{\celsius} 的剩余预算，其值小于仅由实线读出的数值，
因为其他驱动因素同样造成增温。}
\end{omfigure}


\section{海平面}\label{sec:b2:anthropogenic-climate-change:sea-level}

\begin{definition}[海平面收支；热比容海平面上升]\label{def:b2:anthropogenic-climate-change:sea-level}
\emph{海平面收支}\index{海平面收支}是观测到的全球平均海平面变化
与其各项估计贡献之和的比较：增暖海洋的热膨胀、
冰川以及格陵兰和南极冰盖的融化，以及陆地储水的变化。
\emph{热比容海平面上升}\index{热比容海平面上升}是
由海水热膨胀引起、质量保持不变时的那部分海平面上升。
\end{definition}

\begin{proposition}[一层海洋的热膨胀]\label{prop:b2:anthropogenic-climate-change:thermosteric}
厚度为 $H$、热膨胀系数为 $\alpha$ 的一层海洋，
若均匀增温 $\Delta T$、质量与面积保持不变，则升高 $\Delta h = \alpha\,\Delta T\,H$。
\end{proposition}
\begin{proof}
柱体质量 $\rho H A$ 不变（面积 $A$），因此 $\Delta H/H = -\Delta\rho/\rho =
\alpha\,\Delta T$（取一阶近似），其中 $\alpha = -(1/\rho)\,\partial\rho/\partial T$。
\end{proof}
% ledger: teos10:alpha-beta

实际海洋中，$\alpha$ 随温度与压强增大，增温集中在
表层附近。评估报告将结果概括为膨胀效率，即海洋每吸收
\qty{e24}{J} 热量，海平面升高 \qty{0.113}{m}，或整个海洋每升温一度升高 \qty{0.617}{m}
（\emph{中等可信度}）。按其平均深度 \qty{3682}{m}，有效 $\alpha$ 为
$0.617/3682 = \qty{1.68e-4}{\celsius^{-1}}$（见\cref{exo:b2:anthropogenic-climate-change:8}）。
% ledger: ar6ch9:expansion, charette:ocean-volume

下图所示收支在不确定度范围内闭合。1971--2018 年间，各项贡献之和为
\qty{2.00}{mm.\text{年}^{-1}}，观测升高为 \qty{2.33}{mm.\text{年}^{-1}}；2006--2018 年间，贡献之和为
3.61，观测值为 3.69。热膨胀是最大的单项，约占 1971--2018 年总和的一半，但冰的贡献正在追赶：
2006--2018 年间冰川与冰盖合计贡献 \qty{1.62}{mm.\text{年}^{-1}}，超过
热膨胀的贡献。海洋在表层趋于稳定后，深层仍将持续增温，而冰盖
需要数世纪才能作出响应；海平面注定将在数世纪至数千年间继续上升，
并将在数千年内保持高位（\emph{高可信度}）。未来两千年内，
若增温维持在 \qty{1.5}{\celsius}，海平面约升高 2 至 \qty{3}{m}；在 \qty{2}{\celsius} 时升高 2 至
\qty{6}{m}，在 \qty{5}{\celsius} 时升高 19 至 \qty{22}{m}（\emph{低可信度}）。
% ledger: ar6ch9:table95, ar6spm:sea-commit

\begin{omfigure}
\centering
\begin{tikzpicture}
\begin{axis}[ybar stacked, width=0.72\linewidth, height=6.2cm, ymin=0, ymax=4.3, bar width=26pt,
  xtick=data, symbolic x coords={1971--2018, 1993--2018, 2006--2018}, enlarge x limits=0.25,
  ylabel={上升速率（\unit{mm.\text{年}^{-1}}）}, tick label style={font=\scriptsize},
  label style={font=\small}, ytick pos=left, xtick pos=bottom,
  legend style={at={(1.03,0.5)}, anchor=west, font=\scriptsize, draw=black!30},
  legend cell align=left, reverse legend]
% ledger: ar6ch9:table95
\addplot[fill=omDef!70, draw=omDef!90!black] coordinates {(1971--2018,1.01) (1993--2018,1.31) (2006--2018,1.39)};
\addlegendentry{热膨胀}
\addplot[fill=omDef!30, draw=omDef!70!black] coordinates {(1971--2018,0.44) (1993--2018,0.55) (2006--2018,0.62)};
\addlegendentry{冰川}
\addplot[fill=omThm!55, draw=omThm!80!black] coordinates {(1971--2018,0.25) (1993--2018,0.43) (2006--2018,0.63)};
\addlegendentry{格陵兰}
\addplot[fill=omThm!25, draw=omThm!70!black] coordinates {(1971--2018,0.14) (1993--2018,0.25) (2006--2018,0.37)};
\addlegendentry{南极}
\addplot[fill=black!20, draw=black!50] coordinates {(1971--2018,0.15) (1993--2018,0.31) (2006--2018,0.60)};
\addlegendentry{陆地储水}
\node[fill=black, rotate=45, minimum size=5pt, inner sep=0pt] at (axis cs:1971--2018,2.33) {};
\node[fill=black, rotate=45, minimum size=5pt, inner sep=0pt] at (axis cs:1993--2018,3.25) {};
\node[fill=black, rotate=45, minimum size=5pt, inner sep=0pt] at (axis cs:2006--2018,3.69) {};
\end{axis}
\end{tikzpicture}
\par\smallskip{\small \omterm{def:b2:anthropogenic-climate-change:sea-level}{海平面收支}：全球平均海平面上升的观测贡献，
以截至 2018 年的三个时段速率表示，堆叠显示，并以菱形表示观测速率。
堆叠柱与菱形之间的差距即收支未能闭合的部分。数据取自 IPCC 第六
次评估报告，第一工作组，表 9.5。}
\end{omfigure}

\section*{练习}\label{sec:b2:anthropogenic-climate-change:exercises}

除另有说明外，练习所用数据如下：二氧化碳加倍的强迫 $F_{2\times} =
\qty{3.93}{W.m^{-2}}$（基于 $C_0 = \qty{278}{ppm}$）；斯特藩–玻尔兹曼常量 $\sigma =
\qty{5.670e-8}{W.m^{-2}.K^{-4}}$；地球半径 \qty{6371}{km}；一年
为 \qty{3.156e7}{s}；碳和二氧化碳的摩尔质量分别为 12.011 和 \qty{44.009}{g.mol^{-1}}；
以及本章中的反馈与\omterm{def:b2:anthropogenic-climate-change:sea-level}{海平面收支}。
% ledger: ipcc:erf-2xco2, const:sigma, earth:Rmean, ciaaw:abridged-2024, ar6ch7:table710, ar6ch9:table95

\begin{exercise}[$\star$]\label{exo:b2:anthropogenic-climate-change:1}
二氧化碳加倍的强迫在平流层调整后为 \qty{3.75}{W.m^{-2}}；
对流层随后作如下调整：$-0.60$（温度）、+0.22（水汽）、+0.45（云）和
+0.11（陆面）\unit{W.m^{-2}}。(a) 求\omterm{def:b2:anthropogenic-climate-change:erf}{有效辐射强迫} ERF。(b) 它比前一次评估的
\qty{3.71}{W.m^{-2}} 超出百分之几？(c) 哪一项调整为负，为何变暖的
对流层会削弱强迫？
\end{exercise}
% ledger: ar6ch7:2xco2

\begin{exercise}[$\star$]\label{exo:b2:anthropogenic-climate-change:2}
1971 至 2018 年间气候系统获得 \qty{434.9}{ZJ}，2006 至 2018 年间获得 \qty{152.4}{ZJ}，
均取首末两年年中之间的时段计算。(a) 求地球表面积。
(b) 将每项获得量换算为以 \unit{W.m^{-2}} 计的平均失衡。(c) 海洋吸收了
\qty{91}{\percent}，求 1971--2018 年间这部分对应多少泽焦耳？
\end{exercise}
% ledger: ar6ch7:inventory, earth:Rmean

\begin{exercise}[$\star$]\label{exo:b2:anthropogenic-climate-change:3}
(a) 将图中各反馈的中心值相加，并与净反馈比较。(b) 求
以净反馈值求 ECS，并求仅由\omterm{def:b2:anthropogenic-climate-change:feedbacks}{普朗克响应}给出的增温。(c)
其他反馈将普朗克增温放大多少倍？
\end{exercise}

\begin{exercise}[$\star$]\label{exo:b2:anthropogenic-climate-change:4}
(a) 将 TCRE（每 \qty{1000}{Gt} 二氧化碳 \qty{0.45}{\celsius}）换算为
每 \qty{1000}{PgC} 碳的增温量，其可能范围 0.27 至 0.63 亦作同样换算。(b) 将
1850 至 2019 年排放的 \qty{2390}{Gt} 二氧化碳换算为碳的拍克数。(c)
按 TCRE，这些排放对应多少增温？
\end{exercise}
% ledger: ar6spm:tcre, ar6spm:budget

\begin{exercise}[$\star\star$]\label{exo:b2:anthropogenic-climate-change:5}
(a) 地表发出 \qty{398}{W.m^{-2}}，温室效应留住其中的 \qty{159}{W.m^{-2}}。
求 $g$ 及发射 \qty{398}{W.m^{-2}} 的黑体的表面温度。(b) 求
\omterm{def:b2:anthropogenic-climate-change:feedbacks}{普朗克响应}在 \qty{288}{K} 时的值，以及相同 $g$ 下 \qty{298}{K} 时的值。(c) 仅凭
\omterm{def:b2:anthropogenic-climate-change:feedbacks}{普朗克响应}，升温后的地球会更敏感还是更不敏感？
\end{exercise}
% ledger: ipcc:budget-surface, ar6ch7:planck

\begin{exercise}[$\star\star$]\label{exo:b2:anthropogenic-climate-change:6}
卫星给出\omterm{def:b2:anthropogenic-climate-change:feedbacks}{水汽反馈}为 $1.85 \pm 0.32$，研究给出\omterm{def:b2:anthropogenic-climate-change:feedbacks}{递减率反馈}为
$-0.50 \pm 0.20$ \unit{W.m^{-2}.\celsius^{-1}}（一个标准差）。(a) 求二者之和。
(b) 若二者相互独立，则其和的标准差为多少？(c) 模式
给出的和的标准差仅为 \num{0.10}。这对二者的误差说明了什么？
\end{exercise}
% ledger: ar6ch7:wv-lr

\begin{exercise}[$\star\star$]\label{exo:b2:anthropogenic-climate-change:7}
从 1850--1900 到 2006--2019，$\Delta N = \qty{0.59}{W.m^{-2}}$，$\Delta T = \qty{1.03}{\celsius}$
且 $\Delta F = \qty{2.20}{W.m^{-2}}$。(a) 求 $\lambda$ 及有效敏感度。(b) 反馈
参数预计变得不那么负，增量为 $\lambda' = \qty{0.5}{W.m^{-2}.\celsius^{-1}}$，
随增温格局演变。求敏感度 $-F_{2\times}/(\lambda + \lambda')$。
(c) 取 $\lambda' = 0.1$，重做 (b)。
\end{exercise}
% ledger: ar6ch7:energy-budget, ar6ch7:pattern, ar6ch7:pattern-ecs

\begin{exercise}[$\star\star$]\label{exo:b2:anthropogenic-climate-change:8}
(a) 全海洋每升温一度的上升量为 \qty{0.617}{m}，平均深度
为 \qty{3682}{m}，求有效膨胀系数。(b) 海洋在
1971--2018 年间获得 \qty{396}{ZJ}。取膨胀效率为 \qty{0.113}{m} 每 \qty{e24}{J}，求\omterm{def:b2:anthropogenic-climate-change:sea-level}{热比容
海平面上升}量，并与预算中的 \qty{47.5}{mm} 相比较。(c) 厚 \qty{700}{m} 的一层水体升温
\qty{0.2}{\celsius}，取 $\alpha = \qty{2e-4}{\celsius^{-1}}$，问其使海平面上升多少？
\end{exercise}
% ledger: ar6ch9:expansion, charette:ocean-volume, ar6ch7:inventory, ar6ch9:table95

\begin{exercise}[$\star\star$]\label{exo:b2:anthropogenic-climate-change:9}
某模式指纹在四个十年内为 $x = (0.10, 0.20, 0.30, 0.40)$ \unit{\celsius}；
观测值为 $y = (0.12, 0.18, 0.33, 0.41)$ \unit{\celsius}。(a) 求最小二乘标度因子
$\beta$，模型为 $y = \beta x + \varepsilon$。(b) 其标准误差由无强迫模拟估计，
为 \num{0.08}。该响应是否被探测到？是否与模式相符？(c) 若
$\beta$ 为 0.5 而误差相同，将意味着什么？
\end{exercise}

\begin{exercise}[$\star\star\star$]\label{exo:b2:anthropogenic-climate-change:10}
两个指纹，即温室气体 $g = (0.2, 0.4, 0.6, 0.8)$ 与气溶胶 $a = (-0.1, -0.2,
-0.2, -0.3)$，均以 \unit{\celsius} 计、取四个十年，拟合于观测 $y = (0.1,
0.2, 0.4, 0.5)$。(a) 写出最小二乘拟合 $y = \beta_g g + \beta_a a
+ \varepsilon$ 的正规方程并求解。(b) 求最近十年中\omterm{def:b2:anthropogenic-climate-change:da}{归因}于每个驱动因子的增温。
(c) 两指纹相似时，$\beta_g$ 与 $\beta_a$ 的不确定度为何会增大？
\end{exercise}

\begin{exercise}[$\star\star\star$]\label{exo:b2:anthropogenic-climate-change:11}
用\omterm{def:b2:anthropogenic-climate-change:budget}{剩余碳预算}方法，TCRE 为每
\qty{1000}{Gt} 对应 \qty{0.45}{\celsius}，2015 至 2019 年排放的 \qty{210}{Gt}，地球系统反馈取每度
\qty{26}{Gt}。(a) 求 \qty{2}{\celsius} 的预算，相对
2010--2019 年升温 \qty{0.93}{\celsius}，若其他驱动因子再增加 \qty{0.1}{\celsius}。(b) 其他驱动因子
额外增温多少，才能使预算等于评估值 \qty{1350}{Gt}？该值是否在评估范围内？(c) 证明
每 \qty{0.1}{\celsius} 的非 CO$_2$ 增温使任一预算的改变量相同。
\end{exercise}

\begin{exercise}[$\star\star\star$]\label{exo:b2:anthropogenic-climate-change:12}
净反馈的极可能范围为 $-1.81$ 至 \qty{-0.51}{W.m^{-2}.\celsius^{-1}}，中心值约为
$-1.16$。(a) 求两端及中心处的 ECS。(b) 证明关于中心对称的 $\lambda$ 范围
所对应的 ECS 范围并不对称，并指出哪一侧更长。(c) 为何
评估认为排除较高的 ECS 比排除较低的更困难？
\end{exercise}
% ledger: ar6ch7:table710

\begin{problem}[{周末问题——解读预算}]\label{pb:b2:anthropogenic-climate-change:1}
本题沿评估报告的脉络，由辐射强迫一直讨论到
\qty{1.5}{\celsius} 的剩余预算。请利用各练习的数据，并给出 1750--2019 年的 ERF（\unit{W.m^{-2}}）：
二氧化碳 2.16、甲烷 0.54、一氧化二氮 0.21、卤代气体 0.41、臭氧 0.47、
平流层水汽 0.05、气溶胶–辐射 $-0.22$、气溶胶–云 $-0.84$、土地利用 $-0.20$、
雪上黑碳 0.08、凝结尾迹 0.06；二氧化碳在 2019 年为 \qty{410}{ppm} 及
2025 年为 \qty{427.35}{ppm}；2006--2018 年的增热为 \qty{0.79}{W.m^{-2}}；能量收支变化
量为 $\Delta N = 0.59$、$\Delta F = \qty{2.20}{W.m^{-2}}$、$\Delta T = \qty{1.03}{\celsius}$；TCRE
为每 \qty{1000}{Gt} 二氧化碳 \qty{0.45}{\celsius}（很可能范围为 0.27 至 0.63）；\qty{2390}{Gt}
1850--2019 年间排放；2010--2019 年人为增温 \qty{1.07}{\celsius}；\qty{210}{Gt} 排放于
2015 至 2019 年；非 CO$_2$ 的附加增温 \qty{0.1}{\celsius}，即评估范围 0.1
至 \qty{0.2}{\celsius} 的下限；地球系统反馈为每度 \qty{26}{Gt}。结果取三位有效
数字。
% ledger: ipcc:erf-agents, ipcc:erf-2xco2, data:co2-mlo, ar6ch7:inventory, ar6ch7:energy-budget, ar6spm:tcre, ar6spm:budget, ar6ch5:budget-parts

\noindent\textbf{第一部分——辐射强迫。}
\begin{enumerate}
  \item 求四种充分混合温室气体的强迫之和。
  \item 求气溶胶强迫之和。
  \item 将所有人为项求和，并与评估值 \qty{2.72}{W.m^{-2}} 比较。
  \item 用对数定律核对 2019 年二氧化碳的强迫。
  \item 用同一定律求 2025 年二氧化碳的强迫。
  \item 哪一项的不确定度最大？为何它主导了总量的不确定度？
\end{enumerate}

\noindent\textbf{第二部分——能量收支。}
\begin{enumerate}[resume]
  \item 2006--2018 年间气候系统获得了多少能量？以泽焦计。
  \item 其中有多少进入了海洋？
  \item 由能量收支变化求反馈参数。
  \item 求有效\omterm{def:b2:anthropogenic-climate-change:ecs}{平衡气候敏感度}。
  \item 由同一数据求\omterm{def:b2:anthropogenic-climate-change:tcr}{瞬态气候响应}。
  \item 由图中的净反馈求 ECS，并比较两个数值。
  \item 给出历史估计值较低的物理原因。
\end{enumerate}

\noindent\textbf{第三部分——每吨的增温。}
\begin{enumerate}[resume]
  \item 将 TCRE 换算为每 \qty{1000}{PgC} 碳对应的增温量。
  \item 换算其很可能范围。
  \item 将 1850--2019 年的排放量换算为拍克碳。
  \item 按 TCRE，1850--2019 年的排放对应多少增温？
  \item 与 \qty{1.07}{\celsius} 的人为增温相比较，该比较对其他人为驱动因子的
        净效应说明了什么？
\end{enumerate}

\noindent\textbf{第四部分——\qty{1.5}{\celsius} 的剩余预算。}
\begin{enumerate}[resume]
  \item 相对 2010--2019 年，还剩多少增温？
  \item 扣除非 CO$_2$ 增温，再除以 TCRE：从
        2010--2019 年的中期起，还剩多少预算？
  \item 扣除 2015--2019 年的排放。
  \item 扣除剩余增温对应的地球系统反馈。
  \item 若其他驱动因子改为增加 \qty{0.2}{\celsius}，预算将改变多少？
  \item 综上所述：自
        2020 年初起，按 TCRE 计算，\qty{1.5}{\celsius} 的\omterm{def:b2:anthropogenic-climate-change:budget}{剩余碳预算}为多少？请与评估值 \qty{500}{Gt} 比较。
\end{enumerate}
\end{problem}
